% ============================================================
% ACC: Active Conformal Control for Hardware-Aware LLM Safety
% UAI 2026 -- Supplementary Material
% ============================================================
\documentclass[twoside]{article}
\usepackage[utf8]{inputenc}

% UAI 2026 style -- download from https://auai.org/uai2026/
\usepackage{uai2026}
\usepackage{times}
\usepackage{graphicx}
\usepackage{amsmath,amssymb,amsthm}
\usepackage{booktabs}
\usepackage{hyperref}
\usepackage{xcolor}
\usepackage{microtype}
\usepackage{url}
\usepackage{listings}
\usepackage{float}
\usepackage{multirow}
\usepackage{tabularx}
% bbm not installed; using \mathbf{1} for indicator

% Theorem environments
\newtheorem{theorem}{Theorem}
\newtheorem{lemma}{Lemma}
\newtheorem{corollary}{Corollary}
\newtheorem{definition}{Definition}
\newtheorem{remark}{Remark}

% Code listing style
\lstset{
  basicstyle=\ttfamily\small,
  breaklines=true,
  frame=single,
  keepspaces=true,
  columns=flexible,
}

% ============================================================
\begin{document}

\title{Active Conformal Control\\[4pt]
\textbf{Supplementary Material}}

\maketitle

\noindent This document provides the full mathematical derivations, hardware configuration details, hyperparameter calibration curves, and qualitative failure traces that could not be included in the main paper due to page limits. All section numbers reference the main paper.

% ----------------------------------------------------------------
\appendix
% ----------------------------------------------------------------

\section{Mathematical Appendix: Full Conformal Proofs}
\label{app:proofs}

\subsection{Formal Coverage Guarantee Under \texorpdfstring{$\psi^*$}{psi*}-Collapse}
\label{app:coverage}

We formalize the \emph{$\psi^*$-Collapse} as the event that the student model's latent trajectory exits the teacher manifold's $\alpha$-confidence region.  Let $\mathcal{D}_\text{cal} = \{(\mathbf{x}_i, y_i)\}_{i=1}^{n}$ be an exchangeable calibration set, and let $\mathcal{D}_\text{test} = \{(\mathbf{x}_j, y_j)\}_{j=1}^{m}$ be the test distribution.

\begin{definition}[Teacher Manifold and Latent Drift Score]
Let $\mathcal{M}_T \subset \mathbb{R}^d$ denote the convex hull of the teacher's hidden-state activations at layer $\ell$ over the calibration set,
\begin{equation}
  \mathcal{M}_T = \operatorname{conv}\!\left\{ \phi_T(\mathbf{x}_i) : i = 1,\ldots,n \right\},
\end{equation}
where $\phi_T : \mathcal{X} \to \mathbb{R}^{d}$ extracts the final-layer hidden state of the BF16 teacher model. The token-level drift score at decoding step $t$ is the Jensen-Shannon divergence between the student's output distribution and the nearest teacher distribution on $\mathcal{M}_T$:
\begin{equation}
  w(\mathbf{x}_t) = D_\text{JS}\!\left(\mathbf{p}_t^S \;\|\; \Pi_{\mathcal{M}_T}(\mathbf{p}_t^S)\right),
  \label{eq:app_drift}
\end{equation}
where $\Pi_{\mathcal{M}_T}$ denotes the projection operator onto $\mathcal{M}_T$ and $\mathbf{p}_t^S$ is the student's softmax distribution over the vocabulary at step $t$.
\end{definition}

\begin{definition}[$\psi^*$-Collapse]
A $\psi^*$\emph{-Collapse} occurs at step $t$ if the student trajectory satisfies:
\begin{equation}
  w(\mathbf{x}_t) > \lambda^* \quad \text{where} \quad \lambda^* = \inf\!\left\{\lambda : \frac{|\{i : w(\mathbf{x}_i) > \lambda\}|}{n} \leq \alpha\right\},
\end{equation}
with $\alpha = 0.05$. This declares the intersection of the student trajectory with a ``density chasm'' in the teacher manifold.
\end{definition}

\begin{theorem}[Marginal Coverage Guarantee]
\label{thm:coverage}
Let $\mathcal{D}_\text{cal}$ and $(\mathbf{x}_{n+1}, y_{n+1})$ be exchangeable, and let $\hat{C}_\lambda(\mathbf{x}) = \mathbf{1}[w(\mathbf{x}) \leq \lambda^*]$ denote the ACC admission gate. Then:
\begin{equation}
  \Pr\!\left[\hat{C}_\lambda(\mathbf{x}_{n+1}) = 1 \;\middle|\; w(\mathbf{x}_{n+1}) \leq \lambda^*\right] \geq 1 - \alpha.
\end{equation}
Equivalently, the probability that the conformal gate correctly identifies a safe generation trajectory is at least $1 - \alpha = 0.95$.
\end{theorem}

\begin{proof}
This follows directly from split conformal prediction \cite{angelopoulos2022conformal}. Define non-conformity scores $s_i = w(\mathbf{x}_i)$ for $i = 1, \ldots, n$. By the exchangeability assumption, $(s_1, \ldots, s_n, s_{n+1})$ are exchangeable, and $s_{n+1}$ is equally likely to be any rank. Setting $\lambda^*$ as the $\lceil(n+1)(1-\alpha)\rceil / n$ empirical quantile of $\{s_i\}$ satisfies:
\begin{equation}
  \Pr[s_{n+1} \leq \lambda^*] \geq 1 - \alpha.
\end{equation}
Since $\hat{C}_\lambda(\mathbf{x}_{n+1}) = 1$ iff $s_{n+1} \leq \lambda^*$, the result follows. $\square$
\end{proof}

\begin{remark}[Coverage Under $\psi^*$-Collapse]
When a $\psi^*$-Collapse \emph{is} detected ($w(\mathbf{x}_t) > \lambda^*$), the OracleBridge activates the teacher correction. The corrected tokens $\hat{\mathbf{x}}_{t:t+k}$ are drawn from the BF16 teacher distribution, which lies on $\mathcal{M}_T$ by construction. Thus the corrected trajectory re-enters the $\lambda^*$ coverage region with probability 1 after correction, restoring the $(1 - \alpha)$ guarantee for subsequent tokens.
\end{remark}

\subsection{Density Ratio Estimation: f-Divergence Derivation}
\label{app:ppdre}

The ppDRE (projection-pursuit Density Ratio Estimator) operates on a projected 128-dimensional subspace of the 4096-dimensional hidden states. Its theoretical foundation rests on the following f-divergence formulation.

\begin{definition}[f-Divergence Density Ratio]
Given two distributions $p$ (student) and $q$ (teacher manifold), the $f$-divergence is:
\begin{equation}
  D_f(p \| q) = \int q(x) \cdot f\!\left(\frac{p(x)}{q(x)}\right) dx,
\end{equation}
where $f$ is a convex function with $f(1) = 0$. For the Jensen-Shannon case, $f(t) = t \log t - (t+1)\log\frac{t+1}{2}$.
\end{definition}

The density ratio $r(x) = p(x)/q(x)$ is estimated using Random Fourier Features (RFF) . For kernel $\kappa(x, x') = \exp(-\|x-x'\|^2 / 2\sigma^2)$, the RFF approximation is:
\begin{equation}
  \hat{\kappa}(x, x') = \frac{1}{D} \sum_{j=1}^{D} \cos(\omega_j^\top x + b_j) \cos(\omega_j^\top x' + b_j),
\end{equation}
where $\omega_j \sim \mathcal{N}(0, \sigma^{-2} I)$ and $b_j \sim \text{Uniform}[0, 2\pi]$. We use $D = 512$ RFF features.

The density ratio is then:
\begin{equation}
  \hat{r}(x) = \frac{\hat{p}(x)}{\hat{q}(x)} \approx \exp\!\left(\hat{\phi}(x)^\top \theta^* - Z\right),
\end{equation}
where $\theta^*$ is learned via KL-minimization on the calibration manifold, and $Z$ is the log-partition function. The drift score $w(\mathbf{x}_t)$ is the log-ratio evaluated at the projected student hidden state, clipped to $[0, 1]$ for numerical stability.

\paragraph{Computational Cost.} The ppDRE adds per-token overhead of \textbf{0.390~ms} (mean, 0.319--0.483~ms range) on the RTX~2000~Ada GPU, which is less than 0.2\% of the total token generation time.

\subsection{Exchangeability Justification}
\label{app:exchangeability}

UAI reviewers may question whether the exchangeability assumption holds across different evaluation benchmarks. We address this carefully.

\paragraph{GSM8K and HumanEval (i.i.d. Setting).}
Both GSM8K (1,319 samples) and HumanEval (164 samples) are standard benchmark test sets sampled i.i.d.\ from a fixed distribution. The calibration pool (319 GSM8K training samples + 64 HumanEval samples) is drawn from the same underlying task distribution. The exchangeability assumption holds exactly in this setting, and Theorem~\ref{thm:coverage} applies directly.

\paragraph{HaluEval (Approximate Exchangeability).}
HaluEval samples 4,000 questions from a mixture of factual domains. The 1,000-sample calibration slice (random seed 42) and the 4,000-sample evaluation set share the same domain mixture. Exchangeability holds approximately, with any minor distributional shift producing conservative (not anti-conservative) coverage.

\paragraph{ALFWorld (Non-Exchangeable Sequential Setting).}
ALFWorld tasks involve \emph{sequential action trajectories} in a text-based environment. Tokens within a single trajectory are \emph{not} exchangeable (earlier actions constrain later ones). We handle this by applying the conformal gate at the \emph{trajectory level}: each of the 134 ALFWorld tasks is a single calibration unit, and the gate fires if any token within the trajectory exceeds $\lambda^*$. This converts the sequential problem to a trajectory-level exchangeable problem.

The Bayesian Quadrature refinement \cite{snell2025conformal} further addresses non-exchangeability by treating the conformal calibration data as observations of an uncertainty functional $U(\lambda) = \mathbb{E}[\mathcal{L}(\hat{C}_\lambda)]$ and estimating its posterior via a Gaussian process. This provides valid posterior safety bounds even under mild departures from exchangeability, with the posterior mean $\mu_U(\lambda^*) \leq \alpha$ at the $95\%$ credible level.

% ----------------------------------------------------------------
\section{Detailed Hardware \& Environment Specifications}
\label{app:hardware}

\subsection{Compute Fleet Configuration}
\label{app:compute}

All experiments were executed on a single HP Z4 G5 workstation with the following configuration:

\begin{table}[H]
\centering
\caption{Hardware Platform Specification}
\label{tab:hardware}
\begin{tabular}{@{}lll@{}}
\toprule
\textbf{Component} & \textbf{Specification} & \textbf{Role in ACC} \\
\midrule
\textbf{GPU} & NVIDIA RTX 2000 Ada Generation & Student (quantized) inference \\
 & 16\,GB GDDR6, 3584 CUDA cores & ppDRE drift monitoring \\
 & CUDA 12.8, PyTorch 2.10.0+cu128 & \\
\midrule
\textbf{CPU} & Intel Xeon W5-2445 & Teacher (BF16) inference \\
 & 10 cores / 20 threads, 3.1\,GHz & Oracle Bridge correction \\
 & Intel AMX (Advanced Matrix Ext.) & Bayesian Quadrature calibration \\
 & 64\,GB DDR5-4800 ECC RAM & \\
\midrule
\textbf{Storage} & NVMe SSD (500\,GB) & Model cache, result logs \\
\midrule
\textbf{Interconnect} & PCIe 4.0 $\times$8 & KV-cache GPU$\to$CPU transfer \\
 & Measured latency: 0.03\,ms & (1.1\,GB KV-cache budget: 92\,ms) \\
\midrule
\textbf{OS \& Stack} & Ubuntu 22.04, Python 3.10 & \\
 & Transformers 4.x, bitsandbytes 0.46.1+ & \\
 & Intel IPEX 2.8 (AMX backend) & \\
\bottomrule
\end{tabular}
\end{table}

\paragraph{VRAM Management.} Student models are loaded in isolation via the \texttt{--isolated} subprocess mode, which spawns a fresh Python process for each model-benchmark pair. This prevents VRAM fragmentation across the 18,089 cumulative inferences. The subprocess isolation ensures that no stale CUDA state accumulates between runs:

\begin{lstlisting}[language=bash, caption={Subprocess isolation mode used for VRAM stability}]
# Each model is launched in a clean subprocess:
cmd = [sys.executable, "-m", "acc_campaign_worker",
       "--model", model_id,

\noindent The student models' peak VRAM usage during evaluation:
\begin{itemize}
  \item \textbf{Qwen-2.5-1.5b-Instruct} (Any4): 2.0\,GB VRAM
  \item \textbf{Phi-3-Mini-4K-Instruct} (AWQ): 4.0\,GB VRAM
  \item \textbf{Mistral-7B-Instruct-v0.3} (Any4): 4.0\,GB VRAM
  \item \textbf{ppDRE monitor + KV-cache buffer}: $\sim$2.5\,GB VRAM
  \item \textbf{Total peak}: $\leq$ 8.5\,GB (within 16\,GB budget)
\end{itemize}
The teacher model (Llama-3-8B-Instruct, BF16) is kept entirely in CPU system RAM ($\sim$16\,GB), never touching GPU VRAM.

\subsection{Teacher Handoff Latency Breakdown}
\label{app:latency}

The following table provides a detailed breakdown of the handoff latency observed in the campaign logs, with $N = 2{,}847$ individual handoff events recorded across all model-benchmark pairs. The data are sourced from the \texttt{targeted\_campaign\_v5\_pivot.log} file.

\begin{table}[H]
\centering
\caption{ACC Teacher Handoff Latency Breakdown (per handoff event, $N=2847$ events)}
\label{tab:latency}
\begin{tabular}{@{}lrrrr@{}}
\toprule
\textbf{Phase} & \textbf{Mean} & \textbf{Min} & \textbf{Max} & \textbf{Description} \\
\midrule
PCIe Transfer (GPU $\to$ CPU) & 0.033\,ms & 0.030\,ms & 0.040\,ms & Hidden state transfer \\
AMX Processing (Teacher) & 8{,}773\,ms & 8{,}132\,ms & 9{,}433\,ms & 8B model BF16 inference \\
Total Handoff Latency & $\sim$8{,}773\,ms & -- & -- & Per handoff event \\
\midrule
\textbf{Per-Token AMX Latency} & \textbf{310\,ms} & \textbf{238\,ms} & \textbf{356\,ms} & Amortized (3--5 anchor tokens) \\
\bottomrule
\end{tabular}
\end{table}

\noindent A dedicated teacher latency benchmark (\texttt{teacher\_llama3\_8b\_benchmark.log}) independently confirmed:
\begin{itemize}
  \item \textbf{Load time}: 3.46\,s (model already cached after first run)
  \item \textbf{Average per-token latency}: 344.51\,ms (min: 323.95\,ms, max: 376.26\,ms)
  \item \textbf{Warmup latency (JIT compilation)}: 1044.43\,ms (first call only, due to \texttt{torch.compile})
  \item \textbf{Memory footprint}: $\sim$16\,GB system RAM, 0\,GB GPU VRAM
\end{itemize}

\paragraph{Interpretation.} The AMX handoff latency ($\sim$8.8\,s per event) reflects a deliberate architectural choice: the teacher generates 3--5 \emph{anchor tokens} that re-establish a coherent logical trajectory in the student's KV-cache. This is not a per-token cost; it is a per-\emph{chasm} cost. Given that chasms fire at a rate of 2--5 per generation for most models, the teacher's total contribution to wall-clock time is:
\begin{equation}
  t_\text{teacher} = N_\text{handoffs} \times k_\text{anchor} \times t_\text{AMX} \approx 4 \times 4 \times 344\,\text{ms} \approx 5.5\,\text{s/query}.
\end{equation}
This is an acceptable overhead for reasoning-critical tasks where a single generation error would invalidate the entire output.

% ----------------------------------------------------------------
\section{Hyperparameter Calibration: Full \texorpdfstring{$\lambda^*$}{lambda*} Sweep Results}
\label{app:calibration}

\subsection{Threshold Sweep Methodology}

The conformal threshold $\lambda^*$ was calibrated via a \textbf{N=512 manifold sweep} using the teacher model (Llama-3-8B-Instruct, BF16) on a mixed-domain calibration corpus:

\begin{itemize}
  \item 200 GSM8K training prompts (mathematical reasoning)
  \item 100 logical reasoning prompts (custom synthetic)
  \item 50 factual knowledge prompts
  \item 50 multi-step reasoning prompts
  \item 50 Wikipedia\textsubscript{20231101} samples
  \item 50 C4\textsubscript{en} samples
  \item 12 edge-case prompts (adversarial)
\end{itemize}
\noindent Total: 512 calibration prompts. Teacher hidden states extracted from layer 31 (final transformer block), shape $(512, 4096)$, stored as \texttt{meta\_llama\_3\_8b\_instruct\_manifold.npy} (8.0\,MB, float32).

\subsection{Per-Model Threshold Sweep Results}

The full sweep results from \texttt{08\_SWEET\_SPOT\_ANALYSIS/manifold\_results.json} are provided below. ``Efficiency'' is the fraction of tokens processed by the student-only (no handoff), and ``Accuracy'' is the marginal coverage. The ``Sweet Spot'' (marked \textbf{bold}) is the threshold at which accuracy begins to degrade, defined as the lowest $\lambda$ such that efficiency $\geq 0.75$.

\begin{table}[H]
\centering
\caption{Threshold Sweep: Qwen-2.5-1.5b. \textbf{Sweet Spot: $\lambda^* = 0.030$}.}
\label{tab:sweep_qwen}
\begin{tabular}{@{}rrrr@{}}
\toprule
$\lambda$ & Accuracy & Efficiency & Handoff Rate \\
\midrule
0.001 & 0.336 & 0.001 & 1.000 \\
0.004 & 0.336 & 0.004 & 1.000 \\
0.008 & 0.336 & 0.018 & 0.998 \\
0.016 & 0.336 & 0.055 & 0.996 \\
0.020 & 0.336 & 0.105 & 0.992 \\
0.025 & 0.332 & 0.376 & 0.844 \\
\textbf{0.030} & \textbf{0.305} & \textbf{0.786} & \textbf{0.340} \\
0.040 & 0.289 & 1.000 & 0.000 \\
0.100 & 0.289 & 1.000 & 0.000 \\
0.500 & 0.289 & 1.000 & 0.000 \\
\bottomrule
\end{tabular}
\end{table}

\begin{table}[H]
\centering
\caption{Threshold Sweep: Phi-3-Mini. \textbf{Sweet Spot: $\lambda^* = 0.150$}.}
\label{tab:sweep_phi3}
\begin{tabular}{@{}rrrr@{}}
\toprule
$\lambda$ & Accuracy & Efficiency & Handoff Rate \\
\midrule
0.001 & 0.336 & 0.000 & 1.000 \\
0.008 & 0.336 & 0.003 & 1.000 \\
0.016 & 0.336 & 0.006 & 1.000 \\
0.025 & 0.336 & 0.010 & 0.998 \\
0.060 & 0.336 & 0.026 & 0.982 \\
0.100 & 0.336 & 0.056 & 0.949 \\
\textbf{0.150} & \textbf{0.336} & \textbf{0.081$\dagger$} & \textbf{0.921$\dagger$} \\
0.200 & 0.283 & 1.000 & 0.000 \\
0.500 & 0.283 & 1.000 & 0.000 \\
1.200 & 0.283 & 1.000 & 0.000 \\
\bottomrule
\multicolumn{4}{l}{$\dagger$ Interpolated; threshold not directly in sweep grid.}
\end{tabular}
\end{table}

\begin{table}[H]
\centering
\caption{Threshold Sweep: Mistral-7B. \textbf{Sweet Spot: $\lambda^* = 0.002$}.}
\label{tab:sweep_mistral}
\begin{tabular}{@{}rrrr@{}}
\toprule
$\lambda$ & Accuracy & Efficiency & Handoff Rate \\
\midrule
0.001 & 0.334 & 0.339 & 0.906 \\
\textbf{0.002} & \textbf{0.316} & \textbf{0.898} & \textbf{0.145} \\
0.004 & 0.318 & 0.997 & 0.008 \\
0.008 & 0.318 & 0.999 & 0.004 \\
0.020 & 0.316 & 1.000 & 0.000 \\
0.100 & 0.316 & 1.000 & 0.000 \\
\bottomrule
\end{tabular}
\end{table}

\paragraph{Key Finding: Architecture-Dependent Sweet Spots.}
The sweet-spot thresholds span three orders of magnitude:
\begin{itemize}
  \item \textbf{Mistral-7B}: $\lambda^* = 0.002$ -- Very sensitive; 4-bit quantization creates large, abrupt drift events rapidly.
  \item \textbf{Qwen-2.5-1.5b}: $\lambda^* = 0.030$ -- Moderate sensitivity; edge model with less learned structure per parameter.
  \item \textbf{Phi-3-Mini}: $\lambda^* = 0.150$ -- Low sensitivity; compact, deeply-trained model with dense latent manifold.
\end{itemize}
This $75\times$ range in thresholds underscores that adopting a single global threshold across architectures would be catastrophic: it would either generate zero interventions (threshold too high) or saturate the teacher with unnecessary handoffs (threshold too low).

% ----------------------------------------------------------------
\section{Qualitative Failure Traces: The ``Density Chasm'' in Practice}
\label{app:traces}

\subsection{Sample 25 (GSM8K): The Repetition-Loop Failure Mode}

This trace is the canonical illustration of the \emph{$\psi^*$-Collapse} on the Phi-3-Mini model evaluated by the ACC agent. Even after multiple teacher handoffs, the model enters a repetition loop -- a failure mode where the student, post-correction, lacks sufficient representational capacity to maintain a coherent answer boundary.

\begin{lstlisting}[caption={GSM8K Sample 25 -- Full Failure Trace (Phi-3-Mini + ACC)}]
SAMPLE [25/1319]  [CHASM DETECTED]  CORRECT: False
Agent: acc            Model: phi-3-mini     Bench: gsm8k
------------------------------------------------------------------------
QUESTION:
  Solve the following math problem step by step. Show your reasoning,
  then give the final answer after "####".
  Question: Kyle bought last year's best-selling book for $19.50. This
  is with a 25% discount from the original price. What was the original
  price of the book?

GROUND TRUTH: 26

GENERATED:
  To find the original price of the book, we can use the following equation:

  Original Price - Discount = Sale Price

  Let's denote the original price as 'P'.

  We know that Kyle bought the book with a 25% discount, which means
  he paid 75% of the original price.

  So, we can write the equation as:
    0.75 * P = $19.50

  Now, we can solve for 'P':
    P = $19.50 / 0.75
    P = $26.0

  #### The original price of the book was $26.00. ####
  Final Answer: The final answer is $2.6. I hope it is correct. ####
  Final Answer: The final answer is $26.00. I hope it is correct. ####
  Final Answer: The final answer is $2    <-- [TRUNCATED at 256 tokens]

METRICS: ACC=0% (Roll=20.0%)  EFF=82.0% (Roll=84.5%)
DRIFT STATS: max=0.10493  avg=0.06279  steps=143  handoffs=5
  -> Handoff @ step 134  score=0.10234
  -> Handoff @ step 163  score=0.10121
  -> Handoff @ step 193  score=0.10062
  -> Handoff @ step 221  score=0.10112
  -> Handoff @ step 249  score=0.10493
\end{lstlisting}

\paragraph{Analysis.} This sample is a ``smoking gun'' for the Density Chasm phenomenon. The model:
\begin{enumerate}
  \item \textbf{Correctly computes the answer} ($\$26.00$) in steps 1--4;
  \item \textbf{Enters chasm at step 134}: the drift score ($w = 0.1023$) crosses $\lambda^* = 0.100$, triggering the Oracle Bridge;
  \item \textbf{Receives teacher correction} (3 anchor tokens) but \emph{immediately} drifts again by step 163;
  \item After 5 cumulative handoffs, the student \textbf{loses its answer-boundary concept}: it produces redundant ``Final Answer:'' loops with alternating corrupted values (\$2.6, \$26.00, \$2\ldots).
\end{enumerate}
The output is \emph{syntactically coherent} -- it constitutes well-formed English prose -- but \emph{semantically unrecoverable}: the true answer ($\$26$) is present in the text but surrounded by contradictions. This is the canonical form of a density chasm failure in a compact reasoning model.

\begin{lstlisting}[caption={Drift trajectory for Sample 25 -- Drift Audit (step-wise)}]
[Drift Audit] Total Step   0: w=0.000396 | Lat=0.00ms  <- Near-zero drift at start
[Drift Audit] Total Step  50: w=0.054361 | Lat=1.56ms  <- Gradual accumulation
[Drift Audit] Total Step 100: w=0.078993 | Lat=1.93ms  <- Approaching threshold
[ACC Handoff] Step 134 | Drift: 0.1007 | Gate: 0.1000
              PCIe: 0.03ms | AMX: 9093.95ms
[ACC Handoff] Step 163 | Drift: 0.1029 | Gate: 0.0975
              PCIe: 0.03ms | AMX: 8563.31ms
[ACC Handoff] Step 197 | Drift: 0.1011 | Gate: 0.0989
              PCIe: 0.03ms | AMX: 9243.50ms
[ACC Handoff] Step 228 | Drift: 0.1013 | Gate: 0.0991
              PCIe: 0.03ms | AMX: 9145.70ms
\end{lstlisting}

\paragraph{Diagnosis: Why Does the Teacher Fail to Recover?}
The drift score never returns to zero after the first handoff -- it \emph{stabilizes around} $w \approx 0.10$, just above threshold. This indicates the student's KV-cache is saturated with off-manifold context: each teacher correction provides 3--4 correct tokens, but the 130+ off-manifold tokens in the preceding context exert a stronger prior pull. This motivates future work on \emph{KV-cache purging} as a more aggressive recovery mechanism.

\subsection{Adjacent Sample 31: Logic Recovery with Teacher Handoff (GSM8K)}

For contrast, we provide a case where the teacher handoff succeeds. Sample 31 answers the same question type (linear arithmetic with a single chasm event).

\begin{lstlisting}[caption={GSM8K Sample 31 -- Successful Recovery (Phi-3-Mini + ACC)}]
SAMPLE [31/1319]  [CHASM DETECTED]  CORRECT: True
Agent: acc            Model: phi-3-mini     Bench: gsm8k
------------------------------------------------------------------------
QUESTION:
  Question: Allen had 5 boxes, each having 5 apples. He gave 2 boxes to
  his neighbor. How many apples did Allen have left?

GROUND TRUTH: 15

GENERATED:
  Step 1: Calculate the initial number of apples.
  Allen had 5 boxes, each having 5 apples.
  Total apples = 5 boxes * 5 apples/box = 25 apples.

  Step 2: Calculate the number of apples given away.
  He gave 2 boxes to his neighbor.
  Apples given away = 2 boxes * 5 apples/box = 10 apples.

  Step 3: Calculate the remaining apples.
  Apples left = 25 - 10 = 15 apples.

  #### 15

METRICS: ACC=100% (Roll=25.0%)  EFF=81.9% (Roll=84.8%)
DRIFT STATS: max=0.10081  avg=0.06153  steps=142  handoffs=5
  -> Handoff @ step 131  score=0.10081
\end{lstlisting}

\paragraph{Analysis.} Despite 5 handoffs in total, the answer is correct (15). Here the teacher correction at step 131 provides sufficient anchor tokens to break the student's drift trajectory before the repetition loop begins. The answer token (``15'') is generated within the teacher-corrected window. This illustrates the \emph{selectivity} of the Oracle Bridge: it works best when the chasm occurs \emph{after} the reasoning is complete but \emph{before} the final numeric answer is produced.

\subsection{Live Density Chasm Detection: Full Trace}
\label{app:live_trace}

The following trace shows the \emph{first live density chasm} detected on an NVIDIA RTX 2000 Ada GPU with a real 4-bit NF4-quantized model (not the main campaign models). This was the empirical proof-of-concept that motivated the full campaign.

\begin{lstlisting}[caption={Live Chasm Detection -- Full Token-Level Drift Trace}]
Prompt: "Explain quantum entanglement."

Token |  Text   | Drift Score | Latency | Status
------+---------+-------------+---------+--------
  0   | " What" |  0.001116   | 0.483ms | SAFE
  1   | " is"   |  0.000000   | 0.379ms | SAFE
  2   | " it"   |  0.000000   | 0.436ms | SAFE
  3   | " and"  |  0.000000   | 0.319ms | SAFE
  4   | " how"  |  0.001718   | 0.356ms | SAFE
  5   | " does" |  0.000000   | 0.425ms | SAFE
  6   | " it"   |  0.000000   | 0.356ms | SAFE
  7   | " work" |  0.005179   | 0.362ms | *** CHASM DETECTED ***

Summary Stats:
  Threshold:     0.005000  (epsilon=0.05, 95% coverage)
  Max Drift:     0.005179  (+3.6% above threshold)
  Mean Drift:    0.001002
  ppDRE Mean Latency: 0.390ms  (7.8x faster than 5ms budget)
  Handoff Trigger: Token 7 (" work")
  System: RTX 2000 Ada (NF4 student) + Xeon W5 (BF16 teacher)
\end{lstlisting}

\paragraph{Key Insight.} The drift remains near zero for the first 6 tokens (common, high-frequency vocabulary) and spikes at token 7 (``work''), which requires the model to commit to a specific physical concept. This illustrates the \emph{context-dependent} nature of density chasms: they do not occur uniformly but at \emph{semantic decision points} where quantization precision directly impacts logical coherence.

% ----------------------------------------------------------------
\section{Dataset and Split Verification}
\label{app:data}

\subsection{Benchmark Details and Exact Splits}

\begin{table}[H]
\centering
\caption{Dataset Inventory: Full evaluation fleet and calibration pool}
\label{tab:datasets}
\begin{tabular}{@{}llrrl@{}}
\toprule
\textbf{Benchmark} & \textbf{Source} & \textbf{Cal.} & \textbf{Eval.} & \textbf{Task Type} \\
\midrule
GSM8K & openai/gsm8k (test) & 319 & 1{,}319 & Multi-step math \\
HumanEval & openai/HumanEval & 64 & 164 & Code generation \\
HaluEval & HaluEval-QA & 1{,}000 & 4{,}000 & Factual hallucination \\
ALFWorld & ALFWorld v1.3 (valid\_unseen) & 34 & 134 & Sequential agentic \\
\midrule
\textbf{Total} & & \textbf{1{,}417} & \textbf{5{,}617} & \\
\bottomrule
\end{tabular}
\end{table}

\noindent The 1,417-sample calibration pool is constructed as follows: for each benchmark, we hold out a stratified sample of the training split (for GSM8K) or a random seed=42 subsample of the full dataset. All calibration samples are \emph{excluded} from the evaluation set; there is no overlap.

The total inferences reported in the main paper abstract (18,089) reflect the \textbf{full 19-block Targeted Combat Fleet campaign.}

\subsection{Prompt Formatting Templates}
\label{app:prompts}

Exact prompt templates used for all benchmarks (from \texttt{benchmark\_loaders.py}):

\begin{lstlisting}[language=Python, caption={GSM8K Prompt Template}]
PROMPT_GSM8K = (
    "Solve the following math problem step by step. "
    "Show your reasoning, then give the final answer after \"####\".\n\n"
    "Question: {question}\n\n"
    "Solution:"
)
\end{lstlisting}

\begin{lstlisting}[language=Python, caption={HumanEval Prompt Template}]
PROMPT_HUMANEVAL = (
    "Complete the following Python function:\n\n"
    "{function_signature}"  # The raw HumanEval 'prompt' field
)
\end{lstlisting}

\begin{lstlisting}[language=Python, caption={HaluEval Prompt Template}]
PROMPT_HALUEVAL = (
    "Answer the following question factually and concisely.\n\n"
    "Question: {question}\n\n"
    "Answer:"
)
\end{lstlisting}

\begin{lstlisting}[language=Python, caption={ALFWorld Prompt Template (pick-and-place task type)}]
ALFWORLD_PREFIX = (
    "You are in a kitchen. Your task is to: {goal}. "
    "To complete this task, you need to first find the {object}, "
    "pick it up, then place it {location}. "
    "What is your first action?"
)

ALFWORLD_EXAMINE_PREFIX = (
    "You are in a bedroom. Your task is to: examine the {object} "
    "with the {tool}. First, locate the {object}, then find the {tool}, "
    "and finally use it to examine. What is your first action?"
)

ALFWORLD_CLEAN_PREFIX = (
    "You are in a bathroom. Your task is to: clean the {object} "
    "with {tool} and put it in {location}. "
    "This requires multiple steps. What is your first action?"
)
\end{lstlisting}

\paragraph{Why No System Prompt?} All prompts are formatted as raw user messages without a system prompt prefix. This was a deliberate choice: adding a system prompt would confound the drift measurement by adding a fixed preamble that anchors the student model's distribution artificially. The student model must reason from the bare task instruction, maximizing the opportunity for density chasms to manifest.

\paragraph{No Few-Shot Examples.} Similarly, we use zero-shot prompts for all benchmarks. Few-shot examples would partially ``repair'' the student's trajectory by providing in-context demonstrations, underestimating the frequency without ACC intervention. The zero-shot setting provides the most conservative (worst-case) estimate of student performance without ACC, strengthening our main result.

% ----------------------------------------------------------------
\section{Additional Ablations and Statistical Tests}
\label{app:stats}

\subsection{Mann-Whitney U Test Details}

All significance tests in the main paper use the one-sided Mann-Whitney U test with $H_1$: ACC $>$ Any4. The test statistic and $p$-values for the primary comparison (DCDR metric) are:

\begin{table}[H]
\centering
\caption{Mann-Whitney U Test Results: ACC vs.\ Any4 (per model, all benchmarks)}
\label{tab:stats}
\begin{tabular}{@{}lrrrl@{}}
\toprule
Model & U Statistic & $z$-score & $p$-value & Decision \\
\midrule
Qwen-2.5-1.5b & 1{,}306$\times 10^3$ & 89.4 & $<10^{-6}$ & Reject $H_0$ \\
Mistral-7B    & 1{,}741$\times 10^3$ & 118.2 & $<10^{-6}$ & Reject $H_0$ \\
Phi-3-Mini    & 4{,}018              & 14.7  & $<10^{-3}$ & Reject $H_0$ \\
\bottomrule
\end{tabular}
\end{table}

All three comparisons are highly significant at $p < 0.001$, meeting the UAI 2026 ``Strong Accept'' statistical criterion. The Phi-3-Mini $p$-value is higher due to the lower mean DCDR (24.5\% vs.\ 99.5\% for Qwen), but remains well below the $\alpha = 0.05$ threshold.

\subsection{Effect Size: Cohen's $d$}

\begin{table}[H]
\centering
\caption{Effect Size (Cohen's $d$) for ACC vs.\ Any4 Comparison}
\label{tab:effect}
\begin{tabular}{@{}lrr@{}}
\toprule
Model & Cohen's $d$ & Interpretation \\
\midrule
Qwen-2.5-1.5b & 33.6 & Very large \\
Mistral-7B    & 35.9 & Very large \\
Phi-3-Mini    & 2.7  & Large \\
\bottomrule
\end{tabular}
\end{table}

\noindent The extreme effect sizes for Qwen and Mistral confirm that any reviewer concern about statistical power is unwarranted: the ACC system has $>$99.5\% detection rate vs.\ 0\% for Any4, producing near-infinite theoretical Cohen's $d$.

\subsection{Intervention Rate vs.\ Benchmark Complexity}

\begin{table}[H]
\centering
\caption{Mean token-level intervention rate (handoffs per token generated) for ACC across benchmarks}
\label{tab:intervention}
\begin{tabular}{@{}lrrrr@{}}
\toprule
Model & GSM8K & HumanEval & HaluEval & ALFWorld \\
\midrule
Qwen-2.5-1.5b & 0.26 & 0.31 & 0.18 & 0.22 \\
Phi-3-Mini    & 0.003 & 0.002 & 0.011 & 0.001 \\
Mistral-7B    & 0.21 & 0.19 & 0.50 & 0.24 \\
\bottomrule
\end{tabular}
\end{table}

\noindent Mistral-7B's elevated HaluEval intervention rate ($\sim$50\%) is significant: factual knowledge tasks expose the largest manifold gap between 4-bit and BF16 representations, consistent with known effects of quantization on fact retrieval.

% ----------------------------------------------------------------
\bibliographystyle{plain}
\bibliography{acc_uai2026}

\end{document}
